\section{模型的建立与求解}

\subsection{问题一的模型建立}

\subsubsection{问题分析}

题目要求，在只考虑外延层和衬底界面的一次反射、透射所产生的干涉条纹的情形下，建立确定外延层厚度的数学模型。


光源以一定入射角入射外延层表面时，一部分光在表面直接反射，记为{\textbf{反射光1}}；另一部分光折射进入外延层，在衬底界面反射一次后，再从外延层表面折射而出，记为{\textbf{反射光2}}。


实际光源具有一定的宽度，并非理想的一束光。如下图所示，反射光1与反射光2可能相互叠加，产生干涉。


\begin{figure}[H]
\centering
\includegraphics[width=0.45\textwidth,keepaspectratio]{图片/图1.png}
\caption*{图1 光源叠加示意图}
\end{figure}

两束光叠加后的反射光强随波数呈明暗交替的振荡，即干涉条纹。条纹由两个因素共同决定：一是光在传播、折射与反射过程中的强度损耗，它决定两束光各自的振幅，从而决定条纹的幅度；二是两束光的相位差，它决定叠加是加强还是减弱，从而决定条纹的位置。


这两个因素分别携带不同的信息，求解可以分为两步：先由条纹的幅度反推两束光的振幅，进而得到外延层折射率；再由条纹的位置计算外延层厚度。下面依次给出这两步的定量关系。

\subsubsection{求解折射率}

为便于叙述，以下公式的下标中，0 表示空气，1 表示外延层，2 表示衬底。


光具有波的性质。为同时刻画光波的振幅与相位，引入复电场振幅（相量）$\dot{E}$。下面分别分析反射光1与反射光2的光强损耗。设光源的初始复电场振幅为$\dot{E}_0$。


{\textbf{对于反射光1}}，它在空气—外延层上表面直接反射，强度有所损耗。设该界面的反射系数为$r_{01}$，则：
\[
\dot{E}_1 = r_{01}\dot{E}_0 \tag{1}
\]
\indent{\textbf{对于反射光2}}，它依次经历以下过程：（一）从空气透射进入外延层；（二）在外延层中传播到衬底界面；（三）在外延层-衬底界面反射；（四）再次穿过外延层；（五）从外延层透射回空气。
设第一次透射系数为$t_{01}$，外延层—衬底界面的反射系数为$r_{12}$，第二次透射系数为$t_{10}$，光在外延层中往返一次的衰减为$A$，光1与光2的相位差为$\phi$。工作波段内外延层与衬底的吸收较弱，各菲涅耳振幅系数均取实数，材料吸收与传播损耗统一由衰减因子$A$描述，则：
\[
\dot{E}_2 = t_{01}r_{12}t_{10}A \dot{E}_0\cdot e^{i\phi} \tag{2}
\]
\indent 两束光叠加后的复电场振幅为：
\[
\dot{E}_3 = \dot{E}_1+\dot{E}_2 \tag{3}
\]
\indent 叠加后的光强比值（反射率）为：
\[
R = \left|\frac{\dot{E}_3}{\dot{E}_0}\right|^2 \tag{4}
\]

把叠加后的反射率展开。记光2的振幅系数 $B=t_{01}r_{12}t_{10}A$，则：

\[
R=\left|r_{01}+B\,e^{i\phi}\right|^2=r_{01}^2+B^2+2r_{01}B\cos\phi \tag{5}
\]

式(5)中只有 $\cos\phi$ 随波数快速振荡，$r_{01}$ 与 $B$ 变化缓慢，在一个干涉条纹周期内可视为常数。因此，反射率的极值只由 $\cos\phi$ 决定：$\cos\phi$ 取 $+1$ 时两束光同相叠加，反射率达到极大；取 $-1$ 时反相叠加，反射率达到极小。由此得到反射率曲线的上、下包络线：

\[
R_{\max}=(|r_{01}|+|B|)^2,\qquad
R_{\min}=(|r_{01}|-|B|)^2 \tag{6}
\]

由上、下包络线反演薄膜光学常数与厚度的方法称为{\textbf{包络线法}}，本文在包络线的构造上参考了文献\upcite{liu2015}。


式(6)两端均为完全平方：反射率正比于合振幅的平方，而在极值处两束光恰好同相或反相，合振幅等于两束光振幅的和或差。对上、下包络开方，便得到两个可以直接相加减的振幅：

\[
\sqrt{R_{\max}}=|r_{01}|+|B|,\qquad
\sqrt{R_{\min}}=\big||r_{01}|-|B|\big| \tag{7}
\]

这一步是包络线法的关键。若直接对 $R_{\max}$、$R_{\min}$ 作和差，所得只是 $2(r_{01}^2+B^2)$ 与 $4|r_{01}||B|$，$|r_{01}|$ 与 $|B|$ 仍以平方和乘积的形式耦合在一起，无法分离；开方之后，二者均以一次形式出现，可以分别解出：

\[
|r_{01}|=\frac{\sqrt{R_{\max}}+\sqrt{R_{\min}}}{2},
\qquad
|B|=\frac{\sqrt{R_{\max}}-\sqrt{R_{\min}}}{2} \tag{8}
\]

两个振幅的物理含义不同：$r_{01}$ 是空气—外延层界面的反射系数，只由 $n_0$、$n$、入射角和偏振状态决定，{\textbf{与厚度 $d$ 无关}}，携带的是折射率信息；$B$ 则集中了透射、界面反射与往返衰减的全部损耗。因此只需保留前者。记空气—外延层界面的反射系数大小为：

\[
u(\sigma)=\frac{\sqrt{R_{\max}(\sigma)}+\sqrt{R_{\min}(\sigma)}}{2} \tag{9}
\]

上述过程如图2所示。图中是一段由两光束模型生成的示意图：上图给出反射率 $R(\sigma)$ 的振荡曲线及其上、下包络线，并在某一波数处标出这一对极值；下图把两条包络开方。可见 $\sqrt{R_{\max}}$ 与 $\sqrt{R_{\min}}$ 关于 $u(\sigma)$ 对称分布，二者之和的一半即界面反射振幅 $|r_{01}|$，之差的一半即光2的振幅 $|B|$。这样，条纹的幅度信息经过以上三步，就被分解为两束光各自的振幅。


\begin{figure}[H]
\centering
\includegraphics[width=0.85\textwidth,keepaspectratio]{图片/图2 包络线法（由干涉条纹的幅度提取界面反射信息）.png}
\caption*{图2 包络线法（由干涉条纹的幅度提取界面反射信息）}
\end{figure}

得到界面反射系数 $u(\sigma)$ 后，利用斜入射的界面反射公式反演折射率。外延层中的折射角由折射定律确定，$\sin\theta_1(\sigma)=\frac{n_0\sin\theta}{n(\sigma)}$，即：

\[
n(\sigma)\cos\theta_1(\sigma)=\sqrt{n^2(\sigma)-n_0^2\sin^2\theta} \tag{10}
\]

若入射光为 $s$ 偏振，斜入射时的振幅反射系数为：

\[
r_{01}^{(s)}(\sigma)=
\frac{n_0\cos\theta-n(\sigma)\cos\theta_1(\sigma)}
{n_0\cos\theta+n(\sigma)\cos\theta_1(\sigma)} \tag{11}
\]

若入射光为 $p$ 偏振，振幅反射系数为：

\[
r_{01}^{(p)}(\sigma)=
\frac{n(\sigma)\cos\theta-n_0\cos\theta_1(\sigma)}
{n(\sigma)\cos\theta+n_0\cos\theta_1(\sigma)} \tag{12}
\]

接下来处理偏振。由式(11)(12)可见，只要入射角不为零，$s$、$p$ 两种偏振的反射系数就不相等；而题目中的光源不区分偏振，探测器接收到的是两种偏振光强之和。若严格处理，需对两种偏振分别建立式(5)那样的双光束模型，分别计算反射率后相加。这样，实测峰谷包络中混入了两种偏振的贡献，无法由 $u(\sigma)$ 直接解出单一的 $|r_{01}|$，折射率也就没有显式解。对此给出两种处理方案。


{\textbf{方案一：平均近似。}} 把式(11)(12)代入式(5)，记两种偏振各自的反射率为 $R_s=|r_s|^2+B^2-2|r_s|B\cos\phi$ 与 $R_p=|r_p|^2+B^2-2|r_p|B\cos\phi$（干涉项取负，是因为第二束光在衬底界面反射时存在与界面反射系数符号有关的额外相位，使反射率极大出现在 $\cos\phi=-1$ 处；该符号约定与 5.2.4 的实测全谱拟合一致）。对任意入射角 $\theta$，非偏振光的反射率取二者的算术平均：

\[
R(\sigma,\theta)=\frac{R_s+R_p}{2} \tag{13}
\]

{\textbf{方案二：垂直入射近似。}} 当入射角较小时，$\cos\theta\approx\cos\theta_1\approx1$，代入式(11)(12)可知，两种偏振的反射系数退化为同一形式，即垂直入射的菲涅耳公式：

\[
|r_{01}|=\frac{n(\sigma)-n_0}{n(\sigma)+n_0} \tag{14}
\]

由此可以直接解出折射率的显式表达式：

\[
n(\sigma)=n_0\,\frac{1+u(\sigma)}{1-u(\sigma)} \tag{15}
\]

两种方案的比较与选取在 5.2 节进行。

\subsubsection{确定外延层厚度}

上面已由峰谷包络线得到外延层的折射率 $n(\sigma)$。下面利用干涉条纹的{\textbf{位置}}信息，即相邻波峰、波谷的波数，确定厚度 $d$。


\begin{figure}[H]
\centering
\includegraphics[width=0.45\textwidth,keepaspectratio]{图片/图3 光路示意图.png}
\caption*{图3 光路示意图}
\end{figure}

如上图所示，反射光2比反射光1多在外延层中往返一次，光程差为 $2d\,n(\sigma)\cos\theta_1$，对应的相位差为：

\[
\phi=4 \pi \sigma d n(\sigma) \cos \theta_1 \tag{16}
\]

由式(5)可知，波峰处 $\cos\phi=+1$、波谷处 $\cos\phi=-1$，峰、谷对应的相位差分别为 $2m\pi$ 与 $(2m+1)\pi$。因此，从一个波峰到相邻波谷，两束光的相位差改变半个周期：

\[
4\pi d \left[ \sigma_v n(\sigma_v)\cos\theta_{1,v} - \sigma_p n(\sigma_p)\cos\theta_{1,p} \right] = \pi \tag{17}
\]

其中，$p$、$v$ 分别代表波峰、波谷对应的参数。整理可得：

\[
d = \frac{1}{ 4\left| \sigma_v n(\sigma_v)\cos\theta_{1,v} - \sigma_p n(\sigma_p)\cos\theta_{1,p} \right| } \tag{18}
\]

再代入式(10)消去折射角，得到{\textbf{厚度与折射率的函数关系}}：

\[
d= \frac{1}{ 4\left| \sigma_v \sqrt{ n^2(\sigma_v)-n_0^2\sin^2\theta } - \sigma_p \sqrt{ n^2(\sigma_p)-n_0^2\sin^2\theta } \right| } \tag{19}
\]

式(19)中，$\sigma_p$、$\sigma_v$ 由实测反射谱的峰谷位置读出，$\theta$ 由题目给出，$n(\sigma)$ 已由 5.1.2 求得，外延层厚度 $d$ 即可算出。

\subsubsection{总结}

至此，确定外延层厚度所需的公式已全部给出。求解分为两步。


{\textbf{第一步，求折射率。}}
方案1，折射率满足以下方程：
\[
R(\sigma,\theta)=\frac{R_s+R_p}{2} \tag{20}
\]
\indent 其中 $R_s$、$R_p$ 为两种偏振各自的反射率（见 5.1.2）。
方案2，折射率可以显化表示：
\[
n(\sigma)=n_0\,\frac{1+u(\sigma)}{1-u(\sigma)} \tag{21}
\]
\indent 其中，$u(\sigma)=\frac{\sqrt{R_{\max}(\sigma)}+\sqrt{R_{\min}(\sigma)}}{2}$。


{\textbf{第二步，求厚度。}} 把 $n(\sigma)$ 代入厚度公式：
\[
d=\frac{1}{4\left|\sigma_v\sqrt{n^2(\sigma_v)-n_0^2\sin^2\theta}-\sigma_p\sqrt{n^2(\sigma_p)-n_0^2\sin^2\theta}\right|} \tag{22}
\]

其中 $\sigma_p$、$\sigma_v$ 为相邻波峰、波谷的波数，$\theta$ 为入射角，$n_0$ 为空气折射率。

\subsection{问题二的模型建立与求解}

\subsubsection{问题分析}

题目要求，在问题一建立的两光束干涉模型的基础上，设计确定碳化硅外延层厚度的算法，对附件 1（入射角 $10^\circ$）、附件 2（入射角 $15^\circ$）两块晶圆片的实测反射光谱给出厚度计算结果，并分析结果的可靠性。


附件 1、附件 2 是同一晶圆片在两种入射角下的独立实测光谱，数据第 1 列为波数，第 2 列为反射率。由问题一的分析，干涉条纹同时携带两层信息：{\textbf{条纹位置}}反映两束光的相位差，{\textbf{条纹幅度}}反映界面反射强度、只包含折射率 $n$。因此，先从条纹幅度求出折射率 $n(\sigma)$，再把折射率代入条纹位置公式，即可解出厚度 $d$。


算法流程如下图所示。

\begin{figure}[H]
\centering
\includegraphics[width=0.85\textwidth,keepaspectratio]{图片/图4 问题二算法流程图.png}
\caption*{图4 问题二算法流程图}
\end{figure}

其中，第 2、3 两步是核心：第 2 步由条纹幅度反推折射率，第 3 步由条纹位置求出厚度。

\subsubsection{数据预处理}

两份附件的原始反射光谱如下图所示。由图可见，反射率曲线在若干波段出现与干涉条纹正常形态不符的异常特征：部分波数处反射率跳变，谱线起始处还存在一个突跳的零值点。这些异常读数不反映真实的干涉信息，若直接参与后续计算，会污染折射率与厚度的估计。因此，在定量计算之前先对原始数据进行预处理。

\begin{figure}[H]
\centering
\includegraphics[width=0.85\textwidth,keepaspectratio]{图片/图5 原始反射光谱.png}
\caption*{图5 原始反射光谱}
\end{figure}

数据预处理主要包括以下两个方面：


{\textbf{1. 删除反射率为零或负的异常数据点。}} 反射率在物理上不可能取零值或负值，$R\le0$ 的点不可能是任何波段的真实读数。

{\textbf{2. 删除估算折射率超出理论值范围的数据。}} 查阅文献得，碳化硅在红外透明区的折射率范围为 $[2.20,2.60]$\upcite{pitman2008,wang2013}。由相邻极大值与极小值对应的波数反推折射率，去除超出理论范围的峰谷对。


删除异常点后的反射光谱如下图所示，红色三角与绿色倒三角分别标出保留峰谷对的峰、谷位置。

\begin{figure}[H]
\centering
\includegraphics[width=0.85\textwidth,keepaspectratio]{图片/图6 剔除异常点后的反射光谱及干涉峰谷标记.png}
\caption*{图6 剔除异常点后的反射光谱及干涉峰谷标记}
\end{figure}

\subsubsection{方案二的求解}

按问题一的{\textbf{方案二}}（垂直入射近似），折射率可以直接写出解析解。先由相邻峰谷的反射率 $R_{\max}$、$R_{\min}$ 构造上、下包络线：
\[
u(\sigma)=\frac{\sqrt{R_{\max}(\sigma)}+\sqrt{R_{\min}(\sigma)}}{2} \tag{23}
\]
\[
n(\sigma)=n_0\,\frac{1+u(\sigma)}{1-u(\sigma)} \tag{24}
\]

将式(23)(24)应用于处理后的数据，结果列于表 2。由表 2 可见，两份附件所得折射率在数量级与变化趋势上一致，均落在理论范围内，并随波数由 $2.22$ 左右缓慢增大至 $2.48$ 左右。

\begin{longtable}{|>{\centering\arraybackslash}m{0.22\linewidth}|>{\centering\arraybackslash}m{0.2\linewidth}|>{\centering\arraybackslash}m{0.22\linewidth}|>{\centering\arraybackslash}m{0.2\linewidth}|}
\caption{通过理论范围筛选的各峰谷对反推折射率}\\
\hline
中心波数 $\sigma/\mathrm{cm}^{-1}$ & 附件 1（$10^\circ$）$n$ & 中心波数 $\sigma/\mathrm{cm}^{-1}$ & 附件 2（$15^\circ$）$n$ \\[-0.7\baselineskip]
\hline
\endfirsthead
\multicolumn{4}{c}{表\thetable（续）}\\
\hline
中心波数 $\sigma/\mathrm{cm}^{-1}$ & 附件 1（$10^\circ$）$n$ & 中心波数 $\sigma/\mathrm{cm}^{-1}$ & 附件 2（$15^\circ$）$n$ \\
\hline
\endhead
\endfoot
\hline
\endlastfoot
1437.91 & 2.241 & 1353.78 & 2.216 \\
\hline
1534.82 & 2.280 & 1453.34 & 2.305 \\
\hline
1654.38 & 2.327 & 1548.80 & 2.339 \\
\hline
1775.15 & 2.356 & 1671.02 & 2.381 \\
\hline
1892.79 & 2.388 & 1800.71 & 2.411 \\
\hline
2014.28 & 2.403 & 1920.51 & 2.442 \\
\hline
2140.36 & 2.425 & 2032.12 & 2.456 \\
\hline
2261.61 & 2.432 & 2153.62 & 2.477 \\
\hline
2386.72 & 2.447 & — & — \\
\end{longtable}

表 2 中反推折射率止于约 $2.48$，未达到理论范围上界 $2.60$。原因是附件数据在 $2500\ \mathrm{cm}^{-1}$ 以上的干涉条纹幅度过小，未检出峰谷，无法在更高波数处反推折射率。

相邻峰谷对之间的波数间隔较小，折射率在多数波段变化平缓，因此采用{\textbf{线性插值}}构造连续的 $n(\sigma)$ 曲线，供下一步在峰、谷波数处取值。各对反推值之间不作光滑拟合处理，以减少个别波段少量偏差的影响。筛选后的反推折射率随波数的变化如下图所示。

\begin{figure}[H]
\centering
\includegraphics[width=0.85\textwidth,keepaspectratio]{图片/图7 反推折射率随波数的变化.png}
\caption*{图7 反推折射率随波数的变化}
\end{figure}

折射率确定后，由问题一的式(19)求解厚度。每一对相邻峰、谷给出厚度的一个估计值：
\[
d_j=\frac{1}{4\left|\sigma_{v,j}\sqrt{n^2(\sigma_{v,j})-n_0^2\sin^2\theta}-\sigma_{p,j}\sqrt{n^2(\sigma_{p,j})-n_0^2\sin^2\theta}\right|} \tag{25}
\]

式中 $\sigma_{p,j}$、$\sigma_{v,j}$ 分别为第 $j$ 对峰与谷的波数，$n(\sigma_{p,j})$、$n(\sigma_{v,j})$ 由线性插值在相应波数处取值。


$N$ 对峰谷共得到 $N$ 个厚度估计值 $d_j$。由于部分峰谷对反推的折射率可能存在偏差，为减小个别异常值的影响，取{\textbf{中位数作为厚度的最终计算结果}}，并以{\textbf{中位数绝对偏差}}（MAD）描述各估计值的分散程度：
\[
\hat d=\operatorname{median}\{d_j\},\qquad \delta d=\operatorname{median}\{|d_j-\hat d|\} \tag{26}
\]

采用中位数而非平均值，是因为中位数受个别异常值的影响较小。

\subsubsection{方案一的求解}

接下来按问题一的{\textbf{方案一}}（平均近似）计算。两种偏振的界面反射系数取模为 $|r_s|$、$|r_p|$，第二束光振幅为 $B$，相位为 $\phi$。将三者代入 5.1 节所得方程，拟合所用模型为：
{\zihao{5}\[
R(\sigma)=\frac{|r_s(\sigma)|^2+|r_p(\sigma)|^2}{2}+B^2(\sigma)-\big[|r_s(\sigma)|+|r_p(\sigma)|\big]B(\sigma)\cos\left(4\pi\sigma d\sqrt{n(\sigma)^2-n_0^2\sin^2\theta}\right) \tag{27}
\]}

其中 $n(\sigma)$ 由 5.2.3 给出，$B=(\sqrt{R_{\max}}-\sqrt{R_{\min}})/2$ 由 5.1.2 给出，$|r_s|$、$|r_p|$ 由 $n(\sigma)$ 代入 5.1 节的斜入射菲涅耳公式取模得到。因此，上式中唯一待定的量是厚度 $d$。以式(26)所得 $\hat d$ 为初值，在 5.2.2 处理后的有效光谱区内作非线性最小二乘拟合，使模型反射率与实测反射率之差的平方和最小，即可得到 $d$ 的拟合值。

拟合曲线与实测反射率的对比如下图所示，残差曲线绘于右侧纵轴。

\begin{figure}[H]
\centering
\includegraphics[width=0.85\textwidth,keepaspectratio]{图片/图8 全谱拟合结果（附件1，10°）.png}
\caption*{图8 全谱拟合结果（附件1，$10^\circ$）}
\end{figure}

\begin{figure}[H]
\centering
\includegraphics[width=0.85\textwidth,keepaspectratio]{图片/图9 全谱拟合结果（附件2，15°）.png}
\caption*{图9 全谱拟合结果（附件2，$15^\circ$）}
\end{figure}

\subsubsection{计算结果}

{\textbf{方案二的计算结果：}} 按每一对峰谷计算厚度，结果列于表 3。

{\zihao{6}\begin{longtable}{|>{\centering\arraybackslash}m{\dimexpr(\linewidth-2\tabcolsep*(11))/\numexpr10\relax}|>{\centering\arraybackslash}m{\dimexpr(\linewidth-2\tabcolsep*(11))/\numexpr10\relax}|>{\centering\arraybackslash}m{\dimexpr(\linewidth-2\tabcolsep*(11))/\numexpr10\relax}|>{\centering\arraybackslash}m{\dimexpr(\linewidth-2\tabcolsep*(11))/\numexpr10\relax}|>{\centering\arraybackslash}m{\dimexpr(\linewidth-2\tabcolsep*(11))/\numexpr10\relax}|>{\centering\arraybackslash}m{\dimexpr(\linewidth-2\tabcolsep*(11))/\numexpr10\relax}|>{\centering\arraybackslash}m{\dimexpr(\linewidth-2\tabcolsep*(11))/\numexpr10\relax}|>{\centering\arraybackslash}m{\dimexpr(\linewidth-2\tabcolsep*(11))/\numexpr10\relax}|>{\centering\arraybackslash}m{\dimexpr(\linewidth-2\tabcolsep*(11))/\numexpr10\relax}|>{\centering\arraybackslash}m{\dimexpr(\linewidth-2\tabcolsep*(11))/\numexpr10\relax}|}
\caption{通过理论范围筛选的各峰谷对厚度 $d_j$（$\mu\mathrm{m}$）}\\
\hline
\multicolumn{5}{|c|}{附件1（$\theta=10^\circ$）} & \multicolumn{5}{c|}{附件2（$\theta=15^\circ$）} \\
\hline
$\sigma_p$ & $\sigma_v$ & $n(\sigma_p)$ & $n(\sigma_v)$ & $d_j$ & $\sigma_p$ & $\sigma_v$ & $n(\sigma_p)$ & $n(\sigma_v)$ & $d_j$ \\[-0.7\baselineskip]
\hline
\endfirsthead
\multicolumn{10}{c}{表\thetable（续）}\\
\hline
\multicolumn{5}{|c|}{附件1（$\theta=10^\circ$）} & \multicolumn{5}{c|}{附件2（$\theta=15^\circ$）} \\
\hline
$\sigma_p$ & $\sigma_v$ & $n(\sigma_p)$ & $n(\sigma_v)$ & $d_j$ & $\sigma_p$ & $\sigma_v$ & $n(\sigma_p)$ & $n(\sigma_v)$ & $d_j$ \\
\hline
\endhead
\endfoot
\hline
\endlastfoot
1401.51 & 1474.31 & 2.241 & 2.255 & 13.57 & 1291.11 & 1416.46 & 2.216 & 2.272 & 7.03 \\
\hline
1474.31 & 1595.32 & 2.255 & 2.304 & 7.15 & 1416.46 & 1490.22 & 2.272 & 2.318 & 10.60 \\
\hline
1595.32 & 1713.44 & 2.304 & 2.342 & 7.44 & 1490.22 & 1607.38 & 2.318 & 2.359 & 7.43 \\
\hline
1713.44 & 1836.87 & 2.342 & 2.373 & 7.22 & 1607.38 & 1734.66 & 2.359 & 2.396 & 6.90 \\
\hline
1836.87 & 1948.72 & 2.373 & 2.395 & 8.12 & 1734.66 & 1866.76 & 2.396 & 2.428 & 6.65 \\
\hline
1948.72 & 2079.85 & 2.395 & 2.415 & 7.05 & 1866.76 & 1974.27 & 2.428 & 2.449 & 8.31 \\
\hline
2079.85 & 2200.86 & 2.415 & 2.428 & 7.77 & 1974.27 & 2089.98 & 2.449 & 2.466 & 7.88 \\
\hline
2200.86 & 2322.36 & 2.428 & 2.439 & 7.84 & 2089.98 & 2217.26 & 2.466 & 2.477 & 7.45 \\
\hline
2322.36 & 2451.08 & 2.439 & 2.447 & 7.51 & — & — & — & — & — \\
\hline
{\textbf{中位数}} & & & & {\textbf{7.51}} & {\textbf{中位数}} & & & & {\textbf{7.44}} \\
\hline
{\textbf{MAD}} & & & & {\textbf{0.32}} & {\textbf{MAD}} & & & & {\textbf{0.49}} \\
\end{longtable}}

取中位数后得到结果：入射角为 $10^\circ$ 时外延层厚度为 $7.51\ \mu\mathrm{m}$，入射角为 $15^\circ$ 时为 $7.44\ \mu\mathrm{m}$。


{\textbf{方案一的计算结果：}} 入射角为 $10^\circ$ 与 $15^\circ$ 时拟合所得厚度分别为 $7.50\ \mu\mathrm{m}$、$7.27\ \mu\mathrm{m}$，残差均方根约 $0.16\%$、$0.20\%$。

{\textbf{两方案的比较：}} 方案二给出厚度 $7.51\ \mu\mathrm{m}$、$7.44\ \mu\mathrm{m}$，方案一给出 $7.50\ \mu\mathrm{m}$、$7.27\ \mu\mathrm{m}$，两方案的差距分别为 $0.1\%$、$2.2\%$，差距较小。方案一并非完全独立，它以方案二的 $\hat d$ 为初值作迭代拟合，本质上是对方案二的核对。因此，取{\textbf{方案二两中位数的平均作为外延层厚度 $d\approx 7.48\ \mu\mathrm{m}$}}。

\subsubsection{结果可靠性分析}

首先，反推得到的折射率与文献中的变化规律基本一致。表 2 显示，两组数据的折射率均随波数增大而缓慢增大，数值范围为 2.216 至 2.477，与碳化硅在红外透明区的折射率范围相符，说明所选波段中的干涉条纹具有合理的物理基础。

其次，两种入射角得到的厚度中心值较为接近。按相邻峰谷法计算，入射角为 $10^\circ$ 时的厚度中位数为 $7.51\ \mu\mathrm{m}$，入射角为 $15^\circ$ 时为 $7.44\ \mu\mathrm{m}$，两者相差约 1.0\%。入射角改变后，光在外延层中的传播路径会发生变化，而两组结果仍保持接近，说明厚度计算结果具有较好的稳定性。

两种入射角下各峰谷对的厚度估计如下图所示。表 3 中少数峰谷对得到的厚度与多数结果差别较大，说明峰谷位置和条纹幅度仍受局部噪声、峰谷识别和近似条件的影响。因此，本文采用中位数作为最终结果，以减小少数异常峰谷对计算结果的影响；表中的 MAD 用于说明各峰谷对计算结果的离散程度。

\begin{figure}[H]
\centering
\includegraphics[width=0.85\textwidth,keepaspectratio]{图片/图10 双入射角厚度估计对比.png}
\caption*{图10 双入射角厚度估计对比}
\end{figure}

全谱拟合的结果也给出了补充说明。入射角为 $10^\circ$ 时，全谱拟合厚度为 $7.50\ \mu\mathrm{m}$，与相邻峰谷法结果十分接近；入射角为 $15^\circ$ 时，全谱拟合厚度为 $7.27\ \mu\mathrm{m}$，与相邻峰谷法相差 $0.17\ \mu\mathrm{m}$。两种方法给出的厚度量级一致，同时也反映出折射率反演和两光束近似对结果有一定影响。

综合两种入射角的计算结果，本文取相邻峰谷法两个中位数的平均值，得到碳化硅外延层厚度约为 $d=7.48\ \mu\mathrm{m}$。

\subsection{问题三的模型建立与求解}

\subsubsection{问题分析}

题目要求推导产生多光束干涉的必要条件，分析多光束干涉对外延层厚度计算精度可能产生的影响，并确定存在多光束干涉时计算外延层厚度的数学模型。

问题一和问题二中，只保留了两束主要反射光：一束光在外延层上表面直接反射，另一束光进入外延层后，在外延层和衬底的界面反射一次，再从上表面射出。实际测量时，第二束光回到上表面后并不一定全部射出，其中一部分会被反射回外延层。它再次到达下界面后还会发生反射，并重新回到上表面。这样便在原有两束光之外，产生第三束、第四束以及更多反射光。
后续各束光的强度逐次衰减。若衰减较慢，各束光的强度仍不可忽略，且彼此保持稳定的相位关系，它们便会与前两束光叠加，形成多光束干涉，使条纹形状发生改变，进而影响折射率与厚度的计算。

后续求解按照图 11 所示的流程进行，先完成理论推导，再分析实测数据。理论推导分两项：一是把光在外延层内多次往返形成的各束反射光逐一写出并求和，得到多光束干涉的反射率表达式，并推导其产生的必要条件；二是将该式与两光束表达式对比，分析条纹形态的偏差，估计多光束干涉对厚度计算精度的影响。理论部分完成后，先对四份附件逐一判别是否出现多光束干涉，再对判定存在多光束干涉的附件 3 与附件 4 给出硅外延层厚度的计算结果，并用独立的算法加以检验。

\begin{figure}[H]
\centering
\includegraphics[width=0.85\textwidth,keepaspectratio]{图片/图11 问题三求解流程图.png}
\caption*{图11 问题三求解流程图}
\end{figure}

\subsubsection{多光束产生的必要条件}

\begin{figure}[H]
\centering
\includegraphics[width=0.45\textwidth,keepaspectratio]{图片/图12 多光束干涉示意图.png}
\caption*{图12 多光束干涉示意图}
\end{figure}

如上图所示，与只有一次反射、一次透射的两光束情形不同，光在外延层内的多次往返还会产生第三束及以后的反射光。题目所述的多光束干涉，指后续反射光对实测干涉条纹产生不可忽略影响的情形。后续光束的强弱可以用相邻两束出射光的振幅比衡量。下面从多次反射出发，推导这一振幅比，并给出多光束干涉产生的必要条件。

仍沿用问题一的记号，设入射光的复电场振幅为 $\dot E_0$。直接在外延层上表面反射的光为：
\[
\dot E_{\mathrm{a}}=r_{01}\dot E_0 \tag{28}
\]

进入外延层并在下界面反射一次后射出的光为：
\[
\dot E_{\mathrm{b}}=B\dot E_0\mathrm e^{\mathrm i\phi},
\qquad
B=t_{01}r_{12}t_{10}A \tag{29}
\]

其中，$A$ 表示光在外延层内往返一次后的衰减，$\phi$ 为这次往返带来的相位差。式(28)与式(29)就是问题一中的两束反射光。

第二束光回到上表面时，其中一部分会在外延层内侧反射，反射系数为 $r_{10}$。这部分光再到达下界面后又会反射一次，并完成一次新的往返。因此，第三束反射光与第二束光相比，额外多乘了系数 $q=r_{10}r_{12}A\mathrm e^{\mathrm i\phi}$。同理，第四束光比第三束光再多乘一次 $q$，之后各束依次类推。第二束及以后各束光的振幅构成公比为 $q$ 的等比数列。光在每次反射与透射中都有能量损失，公比的模小于 1，该数列收敛。求和后，总反射电场为：
\[
\frac{\dot E_{\mathrm r}}{\dot E_0}
=r_{01}+\frac{B\mathrm e^{\mathrm i\phi}}
{1-r_{10}r_{12}A\mathrm e^{\mathrm i\phi}} \tag{30}
\]

相应的反射率为：
\[
R_{\mathrm{multi}}(\sigma)=
\left|
r_{01}+\frac{B\mathrm e^{\mathrm i\phi}}
{1-r_{10}r_{12}A\mathrm e^{\mathrm i\phi}}
\right|^2,
\qquad
\phi=4\pi\sigma d\sqrt{n^2(\sigma)-n_0^2\sin^2\theta} \tag{31}
\]

式(31)就是计入多次反射后的反射率表达式。等比数列公比的模决定了后续光束的强弱，记相邻两束出射光的振幅比为：
\[
\rho=|r_{10}r_{12}A| \tag{32}
\]

$\rho$ 就是该等比数列公比的模，它把界面反射、膜内传播与吸收损耗对后续光束的影响集中在一起。第三束光的振幅是第二束光的 $\rho$ 倍，之后每束再乘一次 $\rho$，逐次衰减。若 $\rho\ll1$，后续各束光迅速减弱，多次反射的贡献可以忽略，式(31)退化为问题一的两光束模型。只有 $\rho$ 不可忽略时，后续光束才对总反射场有明显贡献，条纹才表现出多光束干涉的特征。因此，界面反射是否较强、吸收损耗是否较小，都体现在 $\rho$ 的大小上。

各束光之间还需保持相干，才能形成稳定的干涉条纹。是否相干取决于相邻两束光的光程差：光程差不超过光源相干长度时，相位差保持稳定，能够发生干涉；超出这一范围后，相位差随机变化，干涉项会被平均掉。相邻两束光的光程差为：
\[
\Delta L=2nd\cos\theta_1 \tag{33}
\]

其中 $\theta_1$ 为膜内折射角，符号 $\lesssim$ 表示“小于或近似等于”，即光程差至多接近相干长度而没有明显超出；相干性随光程差增大逐渐退化，$L_c$ 只是量级意义上的特征值，故此处不用严格的小于等于。相干条件为 $\Delta L\lesssim L_c$，$L_c$ 为光源的相干长度。

综合以上推导，多光束干涉的产生有两个主要要求。第一，后续反射光不能在少数几次往返后迅速衰减，即式(32)的振幅比 $\rho$ 不应过小。第二，参与叠加的各束光之间保持稳定的相位关系，即式(33)的光程差 $\Delta L$ 不超过光源的相干长度 $L_c$。前一条件决定后续光束是否足够强，后一条件决定它们能否稳定叠加。
相干长度可由光谱仪的分辨率估算。光源的相干长度约为谱线宽度的倒数，即 $L_c\approx1/\Delta\sigma$。附件数据的波数采样间隔为 $\Delta\sigma\approx0.48\ \mathrm{cm}^{-1}$，分辨率不小于采样间隔，按采样间隔估计得 $L_c\approx2.1\times10^4\ \mu\mathrm{m}$。取 5.3.5 所得厚度 $d\approx3.4\ \mu\mathrm{m}$ 代入式(33)，相邻光束的光程差约为 $23\ \mu\mathrm{m}$，只有相干长度的千分之一。即使按常规测量常用的 $4\ \mathrm{cm}^{-1}$ 分辨率估计，相干长度仍比光程差大两个量级以上。因此，相干条件充分满足，相邻光束能够稳定叠加。

由于附件仅提供波数与反射率数据，各界面的反射系数与膜内传播损耗无法直接测得。下面利用两光束模型从实测峰谷反射率估计 $\rho$，并结合条纹形态指标进行判别，见 5.3.3。

\subsubsection{多光束干涉对厚度计算精度的影响分析}

式(31)与问题一、问题二的两光束表达式相比，结构上的差别是多出一个分母因子，多光束干涉对厚度计算的影响即来源于此。本节先分析该因子给条纹形态与峰谷反射率读数带来的改变，再估计折射率与厚度的偏差，最后依据实测条纹的形态给出判别多光束干涉的判据。

两光束情形的反射率为 $R=|r_{01}+B\mathrm e^{\mathrm i\phi}|^2$，即问题一的式(5)。$\rho$ 很小时，式(31)的分母接近 1，退化为式(5)，两光束模型是式(31)的最低阶近似。$\rho$ 不可忽略时，分母随相位周期性变化，相当于把第二束及后续各束光的振幅作周期性放大。多光束干涉对厚度计算的影响，全部来自这一分母。

分母的作用最直接地体现在条纹形态上。分母的模在干涉谷对应的相位附近最小，各束光在该处叠加得更充分，谷变得更尖、更深；在干涉峰对应的相位附近分母最大，峰则变得更宽、更平缓。条纹的调制深度也随之改变，相邻峰谷之间大部分数据点会偏向中线一侧。

为直观展示两种条纹的区别，按硅的折射率与实测峰谷反射率确定式(29)、式(31)中的各系数，数值合成两种条纹，如下图所示。图中 $|r_{01}|=0.548$，硅的折射率取 3.42，$B$ 由实测峰谷反射率 $R_{\max}\approx39\%$、$R_{\min}\approx21\%$ 确定。为清晰展示形态差异，合成时 $\rho$ 取 0.20，大于硅晶圆的实际值 0.05，两条曲线的差异方向与实际相同。

\begin{figure}[H]
\centering
\includegraphics[width=0.85\textwidth,keepaspectratio]{图片/图13 两光束与多光束干涉条纹对比.png}
\caption*{图13 两光束与多光束干涉条纹对比}
\end{figure}

由图可见，两光束曲线是标准的余弦形条纹；多光束曲线的谷明显变尖，峰明显变宽，相邻峰谷之间大部分数据点位于中线上方。中线指该峰谷区间内峰、谷反射率的平均值 $R_{\mathrm{mid}}=(R_{\max}+R_{\min})/2$。一个峰谷区间内位于中线以下的点占比，两光束恰为 $50\%$，多光束则明显更小。上图中该差异经过放大，实测硅光谱的偏离幅度较小，但偏离方向一致。因此，这一占比可以作为定量判别的依据。

因此，判别采用两个互补的判据，一个估计相邻光束的振幅比，一个检验条纹形态。振幅比可用两光束模型反推。第二束反射光的振幅 $B$ 由相邻峰、谷的反射率确定，即问题一的式(6)。把它与式(29)对比，即可反推衬底界面的反射系数，进而得到振幅比的估计值：
\[
\hat{\rho}=|r_{10}|\,\hat{r}_{12},
\qquad
\hat{r}_{12}=\frac{B}{t_{01}t_{10}A}\approx\frac{\sqrt{R_{\max}}-\sqrt{R_{\min}}}{2t_{01}t_{10}} \tag{34}
\]

其中 $r_{10}$、$t_{01}$、$t_{10}$ 按垂直入射近似，由外延层折射率的文献名义值算出；外延层在测量波段内吸收较弱，取 $A\approx1$。两个阈值由数值模拟确定：按硅的界面参数合成含分母因子的反射谱，逐次改变振幅比，考察峰谷反射率读数与反推折射率的偏差何时达到不可忽略的水平。若 $\hat{\rho}\lesssim0.03$，第三束光可以忽略，两光束模型足够；若 $\hat{\rho}\gtrsim0.05$，则应采用多光束模型。

振幅比的估计只用到峰谷反射率，属于量级判断。条纹形态是多光束干涉更直接的证据，可作独立的检验。两光束干涉的反射率随相位按余弦变化，相邻峰、谷之间位于中线以下的数据点应约占一半。多光束干涉谷变尖、峰变宽，该占比偏离 $50\%$。记第 $j$ 对相邻峰、谷之间位于中线以下的数据点占比为 $\eta_j$，对全部峰谷对取平均：
\[
\eta=\frac{1}{J}\sum_{j=1}^{J}\eta_j \tag{35}
\]

$\eta$ 的正常波动范围由数值模拟确定。两光束条纹严格呈余弦形，$\eta$ 应恰为 $50\%$；在合成的两光束谱上叠加与实测水平相当的随机扰动后重复统计，得到的波动范围约为 $50\%\pm2\%$，代表测量扰动单独造成的偏离。若 $\eta$ 落在该范围之内，条纹可视为余弦形，两光束模型成立；若显著偏离，则条纹已带有多光束干涉的形态特征。

\subsubsection{判别多光束干涉求解}

按 5.3.3 的判据对四份附件逐一统计。峰、谷检测的显著度门槛与 5.2 节相同，反射率为零或负的点先行剔除。附件 1、2 的统计中剔除了两光束模型失效的残余射线带 $787\sim968\ \mathrm{cm}^{-1}$。每个峰谷区间按式(34)给出 $\hat{\rho}$ 的逐对估计并取中位数，按式(35)统计 $\eta$，结果列于表 4。
{\zihao{6}\begin{longtable}{|>{\centering\arraybackslash}m{0.22\linewidth}|>{\centering\arraybackslash}m{0.12\linewidth}|>{\centering\arraybackslash}m{0.16\linewidth}|>{\centering\arraybackslash}m{0.14\linewidth}|>{\centering\arraybackslash}m{0.16\linewidth}|}
\caption{四份附件的多光束干涉定量判别结果}\\
\hline
附件 & 峰谷段数 & 低于中线占比 $\eta$ & 振幅比 $\hat{\rho}$ & 判别结论 \\[-0.7\baselineskip]
\hline
\endfirsthead
\multicolumn{5}{c}{表\thetable（续）}\\
\hline
附件 & 峰谷段数 & 低于中线占比 $\eta$ & 振幅比 $\hat{\rho}$ & 判别结论 \\
\hline
\endhead
\endfoot
\hline
\endlastfoot
附件 1（SiC，$10^\circ$） & 23 & 50.5\% & 0.013 & 两光束模型适用 \\
\hline
附件 2（SiC，$15^\circ$） & 21 & 51.2\% & 0.015 & 两光束模型适用 \\
\hline
附件 3（Si，$10^\circ$） & 15 & 43.2\% & 0.047 & 存在多光束干涉 \\
\hline
附件 4（Si，$15^\circ$） & 14 & 43.8\% & 0.040 & 存在多光束干涉 \\
\end{longtable}}

碳化硅晶圆片两份附件的 $\eta$ 均落在 $50\%\pm2\%$ 之内，$\hat{\rho}$ 也小于 0.03，条纹呈余弦形，两光束模型足够精确，问题二的厚度结果不受多光束干涉的明显影响。硅晶圆片两份附件的 $\eta$ 均在 $43\%$ 左右，显著低于 $50\%$，$\hat{\rho}$ 已接近 0.05，峰形指标与振幅比两种证据相互印证，且两种入射角下的判断一致。因此判定附件 3 与附件 4 的测试结果中出现了多光束干涉。

判定存在多光束干涉后，反射率表达式由两光束的式(5)变为式(31)。分母因子在干涉谷附近使谷变尖变深，在干涉峰附近使峰变宽变缓，实测峰谷反射率不再是问题一式(6)所给出的包络值，式(8)的振幅分解也随之失效。因此，包络线法不能再用于反演折射率，折射率应取文献值。厚度的计算见 5.3.5。

\subsubsection{硅外延层厚度的计算}
5.3.4 已判定附件 3、4 存在多光束干涉，包络线法不能再用于反演折射率，折射率取硅的文献名义值 $n=3.42$。数据预处理与 5.2 节相同，删除反射率为零或负的点。先按问题一的式(19)，以常数折射率对全部峰谷对逐对求厚度。其中，波数低于 $1500\ \mathrm{cm}^{-1}$ 的峰谷对给出的厚度约为 $3.8\sim3.9\ \mu\mathrm{m}$，比高波数段的 $3.4\ \mu\mathrm{m}$ 偏大 $10\%\sim15\%$，其余各对则彼此接近。由此可见，名义折射率只适用于高波数段；重掺杂衬底在低波数端的折射率偏离名义值较多，并伴随较强的吸收，常数折射率已不能描述全谱。为把低波数段也纳入计算，方案一把衬底的色散写成式(36)的 Drude 介电函数，对全谱段作拟合；方案二沿用相邻峰谷法，补充界面反射相移的色散修正；统计仍不计入波数低于 $1500\ \mathrm{cm}^{-1}$ 的区段，该区段的谷位于谱段边缘附近，识别不完整。
重掺杂衬底的介电函数取 Drude 形式（波数以 $\mathrm{cm}^{-1}$ 计）：
\[
\tilde{\varepsilon}_2(\sigma)=\varepsilon_\infty-\frac{\sigma_{\mathrm{p}}^2}
{\sigma^2+\mathrm{i}\gamma\sigma} \tag{36}
\]

其中 $\sigma_{\mathrm{p}}$ 为等离子体波数，$\gamma$ 为阻尼波数。衬底与外延层同为硅，高频介电常数取 $\varepsilon_\infty=n^2=11.7$。把式(36)代入菲涅耳公式，得到衬底界面的复反射系数 $r_{12}(\sigma)$，其模与幅角都随波数变化；再代入式(31)，即得以 $d$、$\sigma_{\mathrm{p}}$、$\gamma$ 为参数的全谱反射率模型，在全区间 $400\sim4000\ \mathrm{cm}^{-1}$ 上作全局非线性最小二乘拟合。厚度初值由傅里叶基频锚定，以免拟合被低波数端的局部极小带偏；衬底参数先在多组初值上网格搜索，再全参数精修。附件 3 的拟合厚度为 $3.40\ \mu\mathrm{m}$，等离子体波数 $3594\ \mathrm{cm}^{-1}$，阻尼波数 $426\ \mathrm{cm}^{-1}$，残差均方根 $2.60\%$；附件 4 的拟合厚度为 $3.37\ \mu\mathrm{m}$，等离子体波数 $3675\ \mathrm{cm}^{-1}$，阻尼波数 $354\ \mathrm{cm}^{-1}$，残差均方根 $1.32\%$。拟合曲线见图 14。拟合覆盖了含低波数端在内的全谱段，残差主要集中在低波数端的深谷附近，这是用单一 Drude 项描述衬底的简化所致。拟合参数对应的屏蔽等离子体边缘 $\sigma_{\mathrm{p}}/\sqrt{\varepsilon_\infty}$ 约为 $1050\ \mathrm{cm}^{-1}$；取电导有效质量 $m^*\approx0.26m_{\mathrm{e}}$，折算衬底掺杂浓度约 $4\times10^{19}\ \mathrm{cm}^{-3}$，与重掺杂衬底的量级一致。
\begin{figure}[H]
\centering
\includegraphics[width=0.85\textwidth,keepaspectratio]{图片/图14 方案一拟合结果（附件3，10°）.png}
\caption*{图14 方案一拟合结果（附件3，$10^\circ$）}
\end{figure}

\begin{figure}[H]
\centering
\includegraphics[width=0.85\textwidth,keepaspectratio]{图片/图14 方案一拟合结果（附件4，15°）.png}
\caption*{图14 方案一拟合结果（附件4，$15^\circ$）}
\end{figure}

方案二从峰谷位置出发。两光束模型的峰谷位置仅由 $\phi=k\pi$ 决定；在式(31)中，分子与分母的模都只依赖同一组合 $\cos(\phi+\delta_{12})$，反射率随该组合单调变化，因此条纹的极值位于：
\[
\phi(\sigma)+\delta_{12}(\sigma)=k\pi,
\qquad
\delta_{12}=\arg r_{12} \tag{37}
\]

$\delta_{12}$ 不随波数变化时，式(37)退化为两光束的峰谷条件，即问题一的式(19)。存在衬底色散时，$\delta_{12}$ 随波数明显变化，仍按常数相移处理会带来系统偏差。相邻峰、谷的相位差仍为 $\pi$，逐对解出厚度：
\[
d_j=\frac{\pi-\delta_{12}(\sigma_{2,j})+\delta_{12}(\sigma_{1,j})}
{4\pi c\,(\sigma_{2,j}-\sigma_{1,j})},
\qquad
c=\sqrt{n^2-n_0^2\sin^2\theta} \tag{38}
\]

$\delta_{12}(\sigma)$ 由方案一拟合的 Drude 参数算出。式(38)的分子含有两个波数处的相移之差，相移随波数的变化已包含在内。逐对代入的结果列于表 5，未经相移修正的式(19)结果一并列出以便对照。各对估计值取中位数作为厚度结果，以 MAD 描述分散程度，统计方式同 5.2.3。

{\zihao{6}\begin{longtable}{|>{\centering\arraybackslash}m{\dimexpr(\linewidth-2\tabcolsep*(9))/\numexpr8\relax}|>{\centering\arraybackslash}m{\dimexpr(\linewidth-2\tabcolsep*(9))/\numexpr8\relax}|>{\centering\arraybackslash}m{\dimexpr(\linewidth-2\tabcolsep*(9))/\numexpr8\relax}|>{\centering\arraybackslash}m{\dimexpr(\linewidth-2\tabcolsep*(9))/\numexpr8\relax}|>{\centering\arraybackslash}m{\dimexpr(\linewidth-2\tabcolsep*(9))/\numexpr8\relax}|>{\centering\arraybackslash}m{\dimexpr(\linewidth-2\tabcolsep*(9))/\numexpr8\relax}|>{\centering\arraybackslash}m{\dimexpr(\linewidth-2\tabcolsep*(9))/\numexpr8\relax}|>{\centering\arraybackslash}m{\dimexpr(\linewidth-2\tabcolsep*(9))/\numexpr8\relax}|}
\caption{附件 3、4 各峰谷对的厚度 $d_j$（$\mu\mathrm{m}$）}\\
\hline
\multicolumn{4}{|c|}{附件 3（$\theta=10^\circ$）} & \multicolumn{4}{c|}{附件 4（$\theta=15^\circ$）} \\
\hline
$\sigma_1$ & $\sigma_2$ & $d_j$ 未修正 & $d_j$ 修正 & $\sigma_1$ & $\sigma_2$ & $d_j$ 未修正 & $d_j$ 修正 \\[-0.7\baselineskip]
\hline
\endfirsthead
\multicolumn{8}{c}{表\thetable（续）}\\
\hline
\multicolumn{4}{|c|}{附件 3（$\theta=10^\circ$）} & \multicolumn{4}{c|}{附件 4（$\theta=15^\circ$）} \\
\hline
$\sigma_1$ & $\sigma_2$ & $d_j$ 未修正 & $d_j$ 修正 & $\sigma_1$ & $\sigma_2$ & $d_j$ 未修正 & $d_j$ 修正 \\
\hline
\endhead
\endfoot
\hline
\endlastfoot
1508.06 & 1722.60 & 3.41 & 3.34 & 1519.15 & 1736.10 & 3.38 & 3.31 \\
\hline
1722.60 & 1927.02 & 3.58 & 3.53 & 1736.10 & 1941.97 & 3.56 & 3.52 \\
\hline
1927.02 & 2141.56 & 3.41 & 3.38 & 1941.97 & 2160.85 & 3.35 & 3.32 \\
\hline
2141.56 & 2347.43 & 3.56 & 3.53 & 2160.85 & 2387.44 & 3.24 & 3.21 \\
\hline
2347.43 & 2564.86 & 3.37 & 3.35 & 2387.44 & 2587.04 & 3.67 & 3.66 \\
\hline
2564.86 & 2779.89 & 3.40 & 3.39 & 2587.04 & 2802.54 & 3.40 & 3.39 \\
\hline
2779.89 & 2993.46 & 3.43 & 3.41 & 2802.54 & 3009.37 & 3.54 & 3.53 \\
\hline
2993.46 & 3203.66 & 3.48 & 3.47 & 3009.37 & 3245.13 & 3.11 & 3.10 \\
\hline
3203.66 & 3434.60 & 3.17 & 3.16 & 3245.13 & 3450.03 & 3.58 & 3.57 \\
\hline
{\textbf{中位数}} & & {\textbf{3.41}} & {\textbf{3.39}} & {\textbf{中位数}} & & {\textbf{3.40}} & {\textbf{3.39}} \\
\hline
{\textbf{MAD}} & & {\textbf{0.05}} & {\textbf{0.05}} & {\textbf{MAD}} & & {\textbf{0.16}} & {\textbf{0.15}} \\
\end{longtable}}

方案二的结果还可以用傅里叶变换独立检验。干涉条纹相当于反射率随波数作周期振荡，基频只由厚度决定；多光束干涉只会带来整倍频的谐波，不会改变基频。对波数不低于 $1500\ \mathrm{cm}^{-1}$ 的谱段，先插值到均匀的波数间隔，扣除缓慢变化的趋势后作傅里叶变换，由功率谱的主峰读出基频。图 15 给出两份附件的功率谱，主峰尖锐突出，二次、三次谐波清晰可辨，与 5.3.4 的判别结论一致。由基频算得的厚度为 $3.44\ \mu\mathrm{m}$ 与 $3.45\ \mu\mathrm{m}$，比未修正的峰谷中位数偏高约 $1\%$。偏高的原因是傅里叶基频把条纹相位中的相移色散一并计入；频率读数的最小间隔约合厚度 $0.07\ \mu\mathrm{m}$，也限制了两者的比对精度。相移修正后，两份附件的峰谷中位数均为 $3.39\ \mu\mathrm{m}$，与方案一的拟合厚度相符，三种方法在不确定度范围内相互印证。
\begin{figure}[H]
\centering
\includegraphics[width=0.85\textwidth,keepaspectratio]{图片/图15 反射谱的傅里叶功率谱（附件3，10°）.png}
\caption*{图15 反射谱的傅里叶功率谱（附件3，$10^\circ$）}
\end{figure}

\begin{figure}[H]
\centering
\includegraphics[width=0.85\textwidth,keepaspectratio]{图片/图15 反射谱的傅里叶功率谱（附件4，15°）.png}
\caption*{图15 反射谱的傅里叶功率谱（附件4，$15^\circ$）}
\end{figure}

最后分析结果的可靠性。第一，两种入射角得到的厚度彼此接近。入射角改变后，光在外延层中的传播路径随之变化，两组结果仍保持一致，说明厚度计算具有较好的稳定性。第二，各峰谷对估计值的分散程度有限，附件 3 的 MAD 约为中位数的 $1\%$，附件 4 约为 $4\%$，与峰谷位置的识别误差水平相当。第三，方案一利用全谱线形，方案二利用峰谷位置，傅里叶检验利用频域结构，三者依据相互独立，结果彼此一致。综合各结果，硅外延层厚度约为 $d=3.39\ \mu\mathrm{m}$。

\subsubsection{总结}

本节按题目要求逐项给出问题三的回答。多光束干涉产生的必要条件有两个。{\textbf{第一，式(32)定义的相邻出射光束振幅比 $\rho$ 不可过小}}，否则后续光束迅速衰减，式(31)退化为两光束模型。{\textbf{第二，式(33)的相邻光束光程差 $\Delta L$ 应不超过光源的相干长度 $L_c$}}；按采样间隔估计，光程差只有相干长度的千分之一，该条件充分满足。

多光束干涉对厚度计算精度的影响体现在两处。{\textbf{第一，式(31)多出的分母因子使干涉谷变尖、峰变宽}}，实测峰谷反射率不再是两光束模型的包络值，包络线法反演折射率随之失效，折射率只能取文献名义值。{\textbf{第二，重掺杂衬底的界面反射相移随波数变化}}，峰谷位置偏离两光束的极值条件，若不作修正，峰谷法会残留约百分之一量级的系统偏差。

按 5.3.3 的振幅比估计与中线占比两项判据对四份附件逐一判别。附件 3、4 的低于中线占比 $\eta$ 约为 $43\%$，振幅比估计值 $\hat{\rho}$ 约为 $0.04\sim0.05$，频谱中又见整倍频谐波，三种证据一致，判定{\textbf{两份硅晶圆片的测试结果中出现了多光束干涉}}。附件 1、2 的 $\eta$ 均落在 $50\%\pm2\%$ 之内，$\hat{\rho}$ 小于 0.03，条纹呈余弦形，两光束模型足够精确，不存在需要消除的多光束影响，{\textbf{问题二的碳化硅厚度结果不受影响}}。

针对存在多光束干涉的硅晶圆片，厚度计算以式(31)为基本模型，采用两种算法：方案一把衬底色散写成式(36)的 Drude 介电函数，对全谱段作非线性最小二乘拟合；方案二按式(38)逐对求解计入相移色散修正的相邻峰谷厚度。两种入射角下，方案一给出 $3.40\ \mu\mathrm{m}$ 与 $3.37\ \mu\mathrm{m}$，方案二均给出 $3.39\ \mu\mathrm{m}$，傅里叶检验给出 $3.44\ \mu\mathrm{m}$ 与 $3.45\ \mu\mathrm{m}$。综合各结果，{\textbf{硅外延层厚度约为 $d=3.39\ \mu\mathrm{m}$}}。



